\chapter{Pipeline Assincrono em AWS}
\label{chap:pipeline_assincrono_aws}

\chapterepigraph{``To succeed, planning alone is insufficient. One must improvise as well.''}{Isaac Asimov, \textit{Foundation} (1951).}
{\small\itshape\color{AccentGray}
Asimov colocou na boca de Hari Seldon uma tens\~{a}o que qualquer engenheiro de sistemas reconhece~\cite{asimov1951foundation}: o plano \'{e} necess\'{a}rio, mas a realidade exige improvisa\c{c}\~{a}o. Paulo Freire formulou o mesmo princ\'{i}pio como condi\c{c}\~{a}o ontol\'{o}gica: o ser humano \'{e} inacabado, e sua grandeza est\'{a} em permanecer aberto ao inesperado --- em n\~{a}o pretender fechar a conta de forma definitiva~\cite{freire1996autonomia}. O pipeline ass\'{i}ncrono deste cap\'{i}tulo nasce exatamente dessa abertura: o modelo s\'{i}ncrono assume que a resposta chegar\'{a} em 29 segundos; o pipeline ass\'{i}ncrono assume que ela pode n\~{a}o chegar, e projeta o sistema para sobreviver a essa incerteza.\par}
\vspace{0.6em}
% ---------------------------------------------------------------
\section{Escopo do capitulo}
% ---------------------------------------------------------------

\begin{quote}\itshape
\textbf{Onde estamos na hist\'oria:} loop interno, escala de produ\c{c}\~ao.
O agente funciona localmente. Agora ele precisa sobreviver ao mundo real:
timeouts de API Gateway, picos de carga e lat\^encia imprediz\'ivel do LLM.
Este cap\'itulo substitui o modelo s\'incrono por um pipeline ass\'incrono
baseado em SQS + Lambda --- a espinha dorsal sobre a qual o loop externo operar\'a.
\end{quote}

Este capitulo descreve, como laboratorio completo, a arquitetura
assincrona\index{arquitetura assincrona|textbf} usada para superar o limite de
timeout do API Gateway\index{API Gateway|textbf} e processar consultas RAG com
LLM em tempo potencialmente superior a 29 segundos. O fluxo adotado e:
\texttt{API Gateway $\to$ Lambda Controller $\to$ SQS $\to$ Lambda Worker $\to$ DynamoDB $\to$ Lambda Status}.

% ---------------------------------------------------------------
\section{Objetivos de aprendizagem}
% ---------------------------------------------------------------

Ao concluir este capitulo, o leitor sera capaz de:

\begin{itemize}[leftmargin=*]
  \item Explicar por que o processamento sincrono falha para pipelines longos;
  \item Implementar o padrao de enfileiramento com
        SQS\index{SQS|textbf} e persistencia de status em
        DynamoDB\index{DynamoDB|textbf};
  \item Construir endpoints \texttt{/query} e \texttt{/status/\{request\_id\}} com
        contrato consistente;
  \item Validar escalabilidade horizontal do worker e observabilidade de estados.
\end{itemize}

% ---------------------------------------------------------------
\section{Conceitos teoricos}
% ---------------------------------------------------------------

\subsection{Padrao fire-and-poll}

A estrategia adotada e \textit{fire-and-poll}\index{fire-and-poll|textbf}:
o cliente envia a requisicao, recebe um \texttt{request\_id} imediatamente e
consulta o estado periodicamente ate a conclusao. Essa abordagem desacopla o
tempo de resposta da API do tempo de computacao do pipeline. Em termos
arquiteturais, este fluxo e uma variante do padr\~ao \textit{saga orchestration}
aplicado a agentes e tarefas ass\'incronas~\cite{awsagentic2025}.

\subsection{Estados de processamento}

O registro no DynamoDB utiliza estados tipicos:

\begin{itemize}[leftmargin=*]
  \item \texttt{PENDING}: requisicao recebida e enfileirada;
  \item \texttt{PROCESSING}: worker em execucao;
  \item \texttt{COMPLETED}: resposta pronta;
  \item \texttt{FAILED}: erro irrecuperavel com mensagem para auditoria.
\end{itemize}

% ---------------------------------------------------------------
\section{Implementacao orientada por passos}
% ---------------------------------------------------------------

\subsection{Passo 1 --- Submeter consulta e enfileirar}

O endpoint \texttt{POST /query} recebe payload e delega o trabalho ao SQS,
retornando \texttt{202 Accepted} com identificador unico:

\begin{lstlisting}[language=Python,
  caption={Fluxo simplificado do controller: enfileirar e responder rapido.},
  label={lst:controller_simplificado}]
request_id = str(uuid.uuid4())
put_item_dynamodb(request_id, status="PENDING", payload=body)
enviar_mensagem_sqs(request_id=request_id, payload=body)

return {
    "statusCode": 202,
    "body": json.dumps({"request_id": request_id, "status": "PENDING"})
}
\end{lstlisting}

\subsection{Passo 2 --- Processar no worker}

O worker consome mensagens da fila, executa o pipeline e grava o resultado:

\begin{lstlisting}[language=Python,
  caption={Fluxo simplificado do worker com update de status.},
  label={lst:worker_simplificado}]
for msg in mensagens_sqs:
    request_id = msg["request_id"]
    atualizar_status(request_id, "PROCESSING")
    try:
        resposta = executar_pipeline_rag(msg["payload"])
        salvar_resultado(request_id, status="COMPLETED", result=resposta)
    except Exception as exc:
        salvar_resultado(request_id, status="FAILED", error=str(exc))
\end{lstlisting}

\subsection{Passo 3 --- Consultar status}

O endpoint \texttt{GET /status/\{request\_id\}} consulta DynamoDB e retorna o
estado corrente para a interface web ou cliente programatico.

\begin{lstlisting}[language=bash,
  caption={Consulta de status no laboratorio.},
  label={lst:status_call}]
curl "$API_URL/status/$REQUEST_ID"
\end{lstlisting}

\subsection{Passo 4 --- Teste ponta a ponta de laboratorio}

Sequencia recomendada da aula:

\begin{enumerate}[leftmargin=*]
  \item Enviar \texttt{POST /query} com payload valido;
  \item Armazenar \texttt{request\_id};
  \item Executar polling a cada 2 s;
  \item Registrar transicao de estados ate \texttt{COMPLETED}.
\end{enumerate}

\begin{lstlisting}[language=bash,
  caption={Script minimo de polling com shell.},
  label={lst:polling_shell}]
REQUEST_ID="..."
while true; do
  RESP=$(curl -s "$API_URL/status/$REQUEST_ID")
  echo "$RESP"
  echo "$RESP" | findstr COMPLETED >nul && break
  echo "$RESP" | findstr FAILED >nul && break
  timeout /t 2 >nul
done
\end{lstlisting}

\subsection{Passo 5 --- Resultado completo da aula}

Ao final, o aluno deve observar:

\begin{itemize}[leftmargin=*]
  \item aceitacao imediata da requisicao no controller;
  \item processamento desacoplado no worker;
  \item resultado final consultavel por \texttt{request\_id};
  \item ausencia de timeout no cliente para cargas demoradas.
\end{itemize}

\subsection{Passo 6 --- Evidencia operacional via frontend CloudFront}

Em validacao de producao assistida por frontend web (CloudFront), observou-se
o comportamento esperado do padrao \textit{fire-and-poll}. Para duas
requisicoes reais (modos \texttt{persona} e \texttt{twin}), o fluxo evoluiu de
\texttt{PENDING} para \texttt{PROCESSING} e finalizou em \texttt{COMPLETED},
sem expiracao no cliente.

\begin{itemize}[leftmargin=*]
  \item \textbf{Persona:} \texttt{request\_id}: \nolinkurl{ddbc5210-63fc-47bc-a274-283e60be84bc}. \par
        Estado final \texttt{COMPLETED};
  \item \textbf{Twin:} \texttt{request\_id}: \nolinkurl{e5e348e7-d266-4ae9-9f26-eb0e659a56a7}. \par
        Estado final \texttt{COMPLETED}.
\end{itemize}

Essa evidencia reforca que a arquitetura assincrona permanece estavel quando
acessada por uma interface de navegador, preservando o mesmo contrato de API
utilizado por scripts de laboratorio.

\subsubsection*{Validacao na reconstrucao completa da infraestrutura (2026-05-03)}

Apos a reconstrucao total da infraestrutura por produtos (quinze stacks Terraform
independentes) descrita no Cap\'itulo~\ref{chap:cicd_multiambiente_bancario},
repetiu-se o smoke test no ambiente \texttt{dev} para confirmar que o padr\~ao
\textit{fire-and-poll} permanecia funcional com os novos endpoints provisionados:

\begin{itemize}[leftmargin=*]
  \item \textbf{Endpoint API Gateway:}
        \url{https://t9o3dos21k.execute-api.sa-east-1.amazonaws.com};
  \item \textbf{Modo segmento} \textemdash \texttt{request\_id}:
        \nolinkurl{0e73a7f9-84ee-4359-aff3-54bd019b3410}. \par
        Tempo de enfileiramento: $< 1$\,s; processamento total: $\approx 18$\,s; \par
        \texttt{cluster\_id = 3}, segmento \emph{Massa Est\'{a}vel},
        estado final \texttt{COMPLETED};
  \item Modelo de segmenta\c{c}\~ao carregado de
        \texttt{s3://langchain-agent-artifacts-dev/\allowbreak{}clientes-agente/modelo\_clustering\_slim.pkl}.
\end{itemize}

O resultado confirma que a substitui\c{c}\~ao dos recursos AWS (novas ARNs,
novo ID de distribui\c{c}\~ao CloudFront e novo ID de API Gateway) n\~ao altera
o contrato de interac\c{c}\~ao entre cliente e pipeline: o mesmo padr\~ao de
enfileiramento, polling e retorno de resultado opera de forma id\^entica.

% ---------------------------------------------------------------
\section{Plano de testes do capitulo}
% ---------------------------------------------------------------

\subsection{Teste funcional}

\begin{itemize}[leftmargin=*]
  \item \textbf{TC-01:} \texttt{POST /query} retorna \texttt{202} com
        \texttt{request\_id} valido.
  \item \textbf{TC-02:} status evolui de \texttt{PENDING} para
        \texttt{PROCESSING} e depois \texttt{COMPLETED}.
  \item \textbf{TC-03:} para payload invalido, status final deve ser
        \texttt{FAILED} com mensagem de erro rastreavel.
\end{itemize}

\subsection{Teste de desempenho}

\begin{itemize}[leftmargin=*]
  \item \textbf{TP-01:} tempo de resposta do controller $\leq 1$\,s
        para 100 requisicoes seriadas.
  \item \textbf{TP-02:} throughput do worker com lote de 500 mensagens
        sem perda de consistencia em DynamoDB.
\end{itemize}

\subsection{Teste de qualidade de resposta}

\begin{itemize}[leftmargin=*]
  \item \textbf{TQ-01:} comparar 20 respostas assincronas com baseline
        sincrono local e verificar equivalencia semantica.
  \item \textbf{TQ-02:} validar que o modo de polling nao altera o
        conteudo final da resposta.
\end{itemize}

% ---------------------------------------------------------------
\section{Criterios de aceite}
% ---------------------------------------------------------------

\begin{itemize}[leftmargin=*]
  \item TC-01 a TC-03 aprovados.
  \item TP-01 e TP-02 dentro dos limites acordados para \texttt{dev}.
  \item TQ-01 com equivalencia media $\geq 95\%$ frente ao baseline.
  \item Validacao real de polling no frontend com status final
        \texttt{COMPLETED} em modos distintos de consulta.
\end{itemize}

% ---------------------------------------------------------------
\section{Espaco para expansao iterativa}
% ---------------------------------------------------------------

\begin{itemize}[leftmargin=*]
  \item \textbf{v7.1} --- Trocar polling por WebSocket para notificacao ativa.
  \item \textbf{v7.2} --- Adicionar \textit{dead-letter queue}\index{dead-letter queue} e retentativas.
  \item \textbf{Lacuna} --- Modelar SLA distinto por tipo de consulta e modo de agente.
\end{itemize}
